\documentclass[11pt]{article}
\usepackage{amsmath, amssymb, amsthm, amsfonts}
\usepackage[margin=1in]{geometry}
\usepackage{titlesec}
\usepackage{hyperref}

\theoremstyle{plain}
\newtheorem{theorem}{Theorem}
\newtheorem{proposition}[theorem]{Proposition}
\newtheorem{lemma}[theorem]{Lemma}
\newtheorem{corollary}[theorem]{Corollary}

\theoremstyle{definition}
\newtheorem{definition}[theorem]{Definition}
\newtheorem{remark}[theorem]{Remark}

\title{Obstructions to Manifold Realization for Uniform Lattices with Torsion}
\author{}
\date{}

\begin{document}

\maketitle

\begin{abstract}
We address the question of whether a uniform lattice $\Gamma$ in a real semi-simple Lie group, containing non-trivial torsion, can be the fundamental group of a compact manifold $M$ (without boundary) whose universal cover $\tilde{M}$ is acyclic over the rational numbers $\mathbb{Q}$. We prove that such a manifold cannot exist. The presence of torsion imposes structural constraints on the classifying space for proper actions ($\underline{E}\Gamma$), specifically requiring the fixed-point sets of finite subgroups to have vanishing Euler characteristic. For a lattice in a semi-simple Lie group, these fixed-point sets are contractible (Euler characteristic 1), leading to a contradiction.
\end{abstract}

\section{Introduction}

Let $G$ be a connected real semi-simple Lie group with finite center, and let $\Gamma \subset G$ be a uniform lattice. This implies that $\Gamma$ is a discrete subgroup such that the quotient space $\Gamma \backslash G$ is compact. We assume that $\Gamma$ contains elements of finite order (torsion), specifically 2-torsion.

The question is whether there exists a compact manifold $M$ satisfying two conditions:
\begin{enumerate}
    \item The fundamental group $\pi_1(M)$ is isomorphic to $\Gamma$.
    \item The universal covering space $\tilde{M}$ is acyclic over the rational numbers $\mathbb{Q}$ (i.e., $H_k(\tilde{M}; \mathbb{Q}) = 0$ for $k > 0$ and $H_0(\tilde{M}; \mathbb{Q}) \cong \mathbb{Q}$).
\end{enumerate}

If such a manifold exists, the group $\Gamma$ acts freely, properly discontinuously, and cocompactly on the space $\tilde{M}$. In the context of finiteness properties of groups, this implies that $\Gamma$ belongs to the class $FH(\mathbb{Q})$ (groups of finite homological type over $\mathbb{Q}$), admitting a finite dimensional model for $B\Gamma$ up to rational homology.

We demonstrate that the existence of torsion in $\Gamma$, combined with the geometry of the symmetric space associated with $G$, creates an obstruction that prevents such a realization.

\section{The Fixed-Point Set Obstruction}

The obstruction is derived from the theory of proper group actions and their relation to free actions on acyclic spaces.

\subsection{Classifying Space for Proper Actions}
For any discrete group $\Gamma$, there exists a universal space for proper actions, denoted $\underline{E}\Gamma$. This is a $\Gamma$-CW complex characterized by the property that for any finite subgroup $H \subset \Gamma$, the fixed-point set $(\underline{E}\Gamma)^H$ is contractible.

For a uniform lattice $\Gamma$ in a semi-simple Lie group $G$, a canonical finite-dimensional model for $\underline{E}\Gamma$ is the symmetric space $X = G/K$, where $K$ is a maximal compact subgroup of $G$. The action of $\Gamma$ on $X$ is proper and cocompact.

\subsection{Fowler's Necessary Condition}
Recent work by Fowler \cite{Fowler2012} establishes necessary conditions for a group to act freely and cocompactly on a $\mathbb{Q}$-acyclic complex. This result can be viewed as an extension of classical Smith theory to the rational setting, mediated by the proper classifying space.

\begin{theorem}[Fowler \cite{Fowler2012}]
Suppose a group $\Gamma$ admits a finite-dimensional model for $\underline{E}\Gamma$. If $\Gamma$ acts freely and cocompactly on a finite-dimensional complex $\tilde{M}$ which is acyclic over $\mathbb{Q}$, then for every non-trivial finite subgroup $H \subset \Gamma$, the Euler characteristic of the fixed-point set in the proper classifying space must vanish:
\[
\chi\left((\underline{E}\Gamma)^H\right) = 0.
\]
\end{theorem}

\section{Application to Lattices with Torsion}

We apply this theorem to the specific case of the lattice $\Gamma$.

\begin{enumerate}
    \item \textbf{Existence of Finite Subgroups:} By hypothesis, $\Gamma$ contains 2-torsion. Let $\gamma \in \Gamma$ be an element of order 2, and let $H = \langle \gamma \rangle \cong \mathbb{Z}_2$ be the finite subgroup generated by it.

    \item \textbf{Geometry of Fixed Sets:} We use the symmetric space $X = G/K$ as the model for $\underline{E}\Gamma$. By the Cartan Fixed Point Theorem and the structure of symmetric spaces, the fixed-point set of an isometry group $H$ acting on a symmetric space of non-compact type is a totally geodesic submanifold $X^H$.
    
    \item \textbf{Euler Characteristic:} Since $X$ is a Hadamard manifold (simply connected with non-positive curvature), any totally geodesic submanifold $X^H$ is also a Hadamard manifold and is therefore contractible. Consequently, its Euler characteristic is:
    \[
    \chi(X^H) = 1.
    \]

    \item \textbf{The Contradiction:} If the manifold $M$ existed, Fowler's theorem would require $\chi(X^H) = 0$. However, we have established that $\chi(X^H) = 1$. This contradiction ($1 \neq 0$) proves that $\Gamma$ cannot act freely and cocompactly on a $\mathbb{Q}$-acyclic space.
\end{enumerate}

\section{Secondary Obstruction: Virtual Euler Characteristic}

A related obstruction arises from the virtual Euler characteristic. If $M$ is a compact manifold with $\mathbb{Q}$-acyclic universal cover, the Euler characteristic $\chi(M)$ must equal the rational Euler characteristic of the group, $\chi_{\text{virt}}(\Gamma)$.
\[
\chi(M) = \chi_{\text{virt}}(\Gamma) = \text{Vol}(\Gamma \backslash G) \cdot C_G.
\]
For many lattices with torsion (e.g., triangle groups), $\chi_{\text{virt}}(\Gamma)$ is a non-integer rational number. Since $\chi(M)$ must be an integer, this provides an immediate contradiction for generic cases. The fixed-point obstruction detailed in Section 2 is strictly stronger, as it applies even in cases where $\chi_{\text{virt}}(\Gamma)$ happens to be integral (e.g., when $\dim(G/K)$ is odd).

\section{Conclusion}

It is not possible for a uniform lattice $\Gamma$ with torsion to be the fundamental group of a compact manifold whose universal cover is acyclic over $\mathbb{Q}$. The presence of torsion ensures that the fixed-point sets in the proper classifying space are non-empty contractible manifolds (Euler characteristic 1), violating the vanishing condition required for the existence of a free action on a rational acyclic complex.

\begin{thebibliography}{9}
\bibitem{Fowler2012}
J. Fowler, \emph{Finiteness properties for some rational Poincaré duality groups}, Illinois J. Math. \textbf{56} (2012), no. 2, 281--299.

\bibitem{Brown1974}
K. S. Brown, \emph{Euler characteristics of discrete groups and $G$-spaces}, Invent. Math. \textbf{27} (1974), 229--264.

\bibitem{Helgason1978}
S. Helgason, \emph{Differential Geometry, Lie Groups, and Symmetric Spaces}, Academic Press, New York, 1978.
\end{thebibliography}

\end{document}
